\documentclass[12pt,a4paper]{article}
\usepackage[left=15mm,right=15mm,top=14mm,bottom=19mm,footskip=9mm]{geometry}
\usepackage{iftex}
\ifPDFTeX
  \usepackage[T5]{fontenc}
  \usepackage[utf8]{inputenc}
  \usepackage{mathptmx}
\else
  \usepackage{fontspec}
  \IfFontExistsTF{Times New Roman}{\setmainfont{Times New Roman}}{%
    \IfFontExistsTF{Tinos}{\setmainfont{Tinos}}{\setmainfont{TeX Gyre Termes}}}
\fi
\usepackage{amsmath,amssymb}
\usepackage[most]{tcolorbox}
\usepackage{fancyhdr,lastpage}
\usepackage{enumitem}
\definecolor{sectionpurple}{RGB}{155,42,143}
\definecolor{notebackground}{RGB}{229,240,245}
\definecolor{noteborder}{RGB}{255,55,55}
\setlength{\parindent}{0pt}
\setlength{\parskip}{2pt}
\setlength{\headheight}{15pt}
\setlength{\abovedisplayskip}{4pt}
\setlength{\belowdisplayskip}{4pt}
\setlength{\abovedisplayshortskip}{2pt}
\setlength{\belowdisplayshortskip}{2pt}
\pagestyle{fancy}
\fancyhf{}
\renewcommand{\headrulewidth}{0pt}
\renewcommand{\footrulewidth}{0.4pt}
\fancyfoot[L]{\fontsize{8.5}{10}\selectfont Tài liệu tổng hợp theo TDMXX23 (Tư Duy Mở)}
\fancyfoot[R]{\fontsize{8.5}{10}\selectfont Trang \thepage/\pageref{LastPage} -- Bài giảng 08.2}
\newtcolorbox{notebox}{colback=notebackground,colframe=noteborder,
  boxrule=0.8pt,arc=2pt,left=7pt,right=7pt,top=3pt,bottom=3pt,
  before skip=3pt,after skip=3pt}
\newcommand{\partheading}[1]{\par\vspace{5pt}\noindent\textbf{\textcolor{sectionpurple}{#1}}\par\vspace{1pt}}
\newcommand{\example}[2][TDM32]{\noindent\textbf{\textcolor{sectionpurple}{Ví dụ #2:} \textcolor{noteborder}{[#1]}}\enspace}
\newcommand{\cvec}[1]{\overrightarrow{#1}}
\setlist[enumerate]{leftmargin=*,itemsep=1pt,topsep=1pt,parsep=0pt,partopsep=0pt}
\begin{document}
\fontsize{12}{14.6}\selectfont
\begin{center}
  \textbf{BÀI GIẢNG 08 -- PHẦN 2}\\[-1pt]
  \textbf{PHƯƠNG PHÁP TỌA ĐỘ HÓA TRONG KHÔNG GIAN OXYZ}\\[-2pt]
  {\fontsize{10.5}{12}\selectfont\itshape Biên soạn theo chủ đề TDMXX23 -- Tư Duy Mở}
\end{center}
\vspace{-7pt}

\partheading{PHẦN I. PHƯƠNG PHÁP TỌA ĐỘ HÓA}
\textbf{Định nghĩa:} Phương pháp tọa độ hóa không gian $Oxyz$ là phương pháp
gắn hệ trục tọa độ hợp lí vào một hình không gian. Từ đó, chúng ta áp dụng
những kiến thức giải tích trong không gian $Oxyz$ đã học để giải quyết
bài toán một cách dễ dàng và mạch lạc.

\textbf{Phương pháp:}
\begin{enumerate}[label=\textbf{Bước \arabic*:}]
  \item Phác họa hình không gian rồi gắn hệ trục tọa độ $Oxyz$ vào một
  cách hợp lí. Chuẩn hóa đơn vị nếu có thể, chẳng hạn chọn độ dài cạnh
  hình lập phương (hoặc một cạnh thích hợp) bằng $1$.
  Sau khi chuẩn hóa, cần hiệu chỉnh kết quả:
  khoảng cách $d$ nhân với $a$; diện tích $S$ nhân với $a^2$;
  thể tích $V$ nhân với $a^3$.
  \item Suy ra tọa độ các điểm tương ứng cần thiết.
  \item Dùng các kiến thức giải tích $Oxyz$ để giải quyết các vấn đề
  của bài toán.
\end{enumerate}

\partheading{PHẦN II. VÍ DỤ ÁP DỤNG}
\begin{notebox}
\example[TDM32]{01} Cho hình lập phương $ABCD.A'B'C'D'$ có cạnh bằng $a$.
Gọi $M$ là trung điểm của cạnh $BB'$, điểm $N$ thỏa mãn
$\cvec{NC}+2\cvec{NC'}=\vec{0}$.
Gọi $\varphi=\bigl((MA'D'),(NAB)\bigr)$.
Hãy tính $\cos\varphi$?
\end{notebox}
\textbf{\textcolor{sectionpurple}{Giải:}}

\begin{notebox}
\example[TDM32]{02} Hãy gắn hệ trục tọa độ $Oxyz$ một cách hợp lí vào
các hình sau:
\begin{enumerate}[label=\textbf{\arabic*.}]
  \item Hình lập phương; hình hộp chữ nhật; hình chóp $S.ABCD$
  có $SA\perp(ABCD)$ và $AB\perp AD$; hình lăng trụ đứng
  $ABC.A'B'C'$ có $AB\perp AC$? Các mặt phẳng $(ABC)$ và $(DBC)$?
  \item Hình chóp đều $S.ABC$?
  \item Hình chóp đều $S.ABCD$?
  \item Hình chóp $S.ABC$ có $\triangle ABC$ đều cạnh $a$,
  $SA=2a$ và $SA\perp(ABC)$. Gọi $M$ là trung điểm của $SB$
  và $N$ là điểm thỏa mãn $\cvec{NA}+2\cvec{NC}=\vec{0}$.
  Hãy tính khoảng cách $MN$?
\end{enumerate}
\end{notebox}
\textbf{\textcolor{sectionpurple}{Giải:}}

\begin{center}
\textbf{\textcolor{sectionpurple}{--- Hết ---}}
\end{center}
\end{document}
